% Copied from neurovlm/scripts/figures/review by export_overleaf.py; edit the repo file and export again.
\documentclass[11pt]{article}
% Response to reviewers, NeuroVLM revision.  Response text: RESPONSE.md (same words; check_letter_sync.py compares them);
% comments: reviews.md (verbatim); figures from reviewer_01/ .. reviewer_03/; figure and table captions are written here
% only (build_reviewer_dirs.py copies the figure captions into the reviewer folders).  Compile here, twice (contents,
% figure numbers and comment links): pdflatex response_letter.tex
\usepackage[utf8]{inputenc}
\usepackage[T1]{fontenc}
\usepackage[letterpaper,margin=1in]{geometry}  % text block 6.5 x 9 in; figures are set at the full text width
\usepackage{graphicx,xcolor,booktabs,array,fancyhdr,etoolbox,placeins}
\newcolumntype{L}[1]{>{\raggedright\arraybackslash}p{#1}}  % ragged-right table column
\usepackage{hyperref}
\hypersetup{colorlinks=true,linkcolor=blue!55!black,urlcolor=blue!55!black,linktoc=all,
  pdftitle={NeuroVLM: point-by-point response to the reviewers}}
\setlength{\parindent}{0pt}
\setlength{\parskip}{0.6em}
\setlength{\tabcolsep}{4pt}
\setlength{\headheight}{14pt}
\pagestyle{fancy}
\fancyhf{}
\renewcommand{\headrulewidth}{0pt}
\fancyhead[L]{\small\color{black!55}NeuroVLM: response to the reviewers}
\fancyhead[R]{\small\hyperlink{contents}{Back to contents}}
\fancyfoot[C]{\small\thepage}
% contents on the title page: less space before each reviewer
\makeatletter\patchcmd{\l@section}{1.0em}{0.5em}{}{}\makeatother
% figures float (here, at the top of a page or on a page of their own) but never past the end of their comment (\FloatBarrier)
\renewcommand{\topfraction}{0.9}\renewcommand{\bottomfraction}{0.8}\renewcommand{\textfraction}{0.07}
\renewcommand{\floatpagefraction}{0.75}\setcounter{topnumber}{3}\setcounter{totalnumber}{4}
\setlength{\textfloatsep}{12pt plus 4pt minus 4pt}\setlength{\intextsep}{10pt plus 4pt minus 4pt}
% figures: Fig. R1, R2, ... and tables: Table R1, R2, ..., both in order of use; captions in footnote size
\renewcommand{\figurename}{Fig.}
\renewcommand{\thefigure}{R\arabic{figure}}
\renewcommand{\thetable}{R\arabic{table}}
\makeatletter
\long\def\@makecaption#1#2{\vskip\abovecaptionskip{\footnotesize\textbf{#1.} #2\par}\vskip\belowcaptionskip}
\makeatother
% the reviewer's comment: verbatim, in a shaded box; \revqpar starts a new paragraph (and box), so that a long comment
% can break across pages between paragraphs
\definecolor{quotebg}{gray}{0.93}
\newsavebox{\revbox}
\newcommand{\revqopen}{\begin{lrbox}{\revbox}\begin{minipage}{\dimexpr\linewidth-2\fboxsep\relax}%
  \setlength{\parskip}{0.45em}\small\itshape\color{black!85}}
\newcommand{\revqclose}{\end{minipage}\end{lrbox}\colorbox{quotebg}{\usebox{\revbox}}}
\newcommand{\revqpar}{\revqclose\par\nointerlineskip\noindent\revqopen}
\newenvironment{revquote}{\par\smallskip\setlength{\parskip}{0pt}\noindent\revqopen
  {\small\upshape\bfseries\color{black!70}Reviewer's comment}\par}%
  {\revqclose\par\medskip}
\newcommand{\response}{\textbf{Response:} }
\newcommand{\revchanges}[1]{\par{\small\raggedright\textbf{Changes to the manuscript:} #1\par}}  % ragged: line references do not break
% references to the revised manuscript in the changes lines: \msl{359}{363} prints "l. 359--363" (the margin numbers of
% the marked manuscript, paper/sn-article-diff.pdf, manuscript_changes.md), \msp{14} "p. 14", \mspp{11}{13} "pp. 11--12".  Plain text here;
% response_letter_with_manuscript.tex defines \mslink, which links each to the appended manuscript page that holds it
% (the paper's line ruler has 46 numbers per page, so line l is on page (l + 22) / 46 as \numexpr rounds it); each
% reference is one unbreakable box, so that a link never spans two lines or pages
\providecommand{\mslink}[2]{#2}
\newcommand{\msl}[2]{\begingroup\edef\x{\endgroup\noexpand\mslink{\number\numexpr(#1+22)/46\relax}}%
  \x{\mbox{l.~#1\ifnum#1=#2\else--#2\fi}}}
\newcommand{\msp}[1]{\mslink{#1}{p.~#1}}
\newcommand{\mspp}[2]{\mslink{#1}{\mbox{pp.~#1--#2}}}
% manuscript figures and tables: \msfig{S8} prints "Fig.~S8", \mstab{S5} "Table~S5", and \msfign{S4} or \mstabn{S6}
% the number alone, after "Figs." or "Tables".  Plain text here; response_letter_with_manuscript.tex defines \mspg, a
% table of the marked page that holds each caption (written by mspages.py), and links each mention to that page
\providecommand{\mspg}[1]{0}
\DeclareRobustCommand{\mscap}[2]{\begingroup\edef\x{\endgroup\noexpand\mslink{\mspg{#1}}}\x{#2}}  % robust: also in captions
\newcommand{\msfig}[1]{\mscap{Fig.#1}{Fig.~#1}}
\newcommand{\mstab}[1]{\mscap{Table.#1}{Table~#1}}
\newcommand{\msfign}[1]{\mscap{Fig.#1}{#1}}
\newcommand{\mstabn}[1]{\mscap{Table.#1}{#1}}
% \authornote{text}: a fact only the authors can supply, printed red on yellow
\newcommand{\authornote}[1]{\fcolorbox{red}{yellow}{\parbox{0.9\linewidth}{\centering\small\color{red}\textbf{#1}}}}
% \cmt{2.4}: link to the heading of comment 2.4
\newcommand{\cmt}[1]{\hyperref[sec:#1]{#1}}
% \rfig{label}{file}{caption}: a numbered figure at the full text width (narrower only if it would not fit the page), or
% at the fraction of it set by \rfigwidth{label}{fraction}: paper figures at their print width (5.15 in), letter plots
% (drawn 6.5 in wide with 8 pt text) at the width their content needs
\newcommand{\rfigwidth}[2]{\expandafter\def\csname rfigw@#1\endcsname{#2}}
\newcommand{\rfig}[3]{\begin{figure}[htp]\centering
  \includegraphics[width=\ifcsname rfigw@#1\endcsname\csname rfigw@#1\endcsname\fi\textwidth,height=0.8\textheight,keepaspectratio]{#2}
  \caption{#3}\label{#1}\end{figure}}
\rfigwidth{fig:fig4}{0.79}\rfigwidth{fig:fig2}{0.79}\rfigwidth{fig:fig3}{0.79}\rfigwidth{fig:figS4}{0.79}
\rfigwidth{fig:fig5}{0.79}\rfigwidth{fig:figS3}{0.5}\rfigwidth{fig:bounds}{0.8}\rfigwidth{fig:ood}{1}
\rfigwidth{fig:mlp}{0.85}\rfigwidth{fig:granularity}{0.85}\rfigwidth{fig:encoder}{0.8}\rfigwidth{fig:latent}{0.85}
\rfigwidth{fig:probes}{0.85}\rfigwidth{fig:steering}{0.85}
\newcommand{\fig}[3]{\rfig{}{#2}{#3}}  % old form, unlabeled
% \rfigparts{label}{file without .png}{caption}: a figure with several pages (or too tall for one), one letter page each,
% as listed by build_reviewer_dirs.py in <file>_parts.tex; a plain \rfig if there is no such file
\newcommand{\rfiglabel}{}\newcommand{\rfigcaption}{}
\newcommand{\rfigparts}[3]{\IfFileExists{#2_parts.tex}{\renewcommand{\rfiglabel}{#1}\renewcommand{\rfigcaption}{#3}%
  \input{#2_parts.tex}}{\rfig{#1}{#2.png}{#3}}}
\newcommand{\rfigpart}[3]{\begin{figure}[htp]\centering
  \includegraphics[width=\textwidth,height=0.84\textheight,keepaspectratio]{#1}
  \ifnum#2=1 \caption{\rfigcaption{} (Page 1 of #3.)}\expandafter\label\expandafter{\rfiglabel}%
  \else\par\small Fig.~\ref{\rfiglabel}, page #2 of #3.\fi\end{figure}}
% \rtab{caption}{tabular}: a numbered table where it stands, caption above, kept on one page with its table
\newcommand{\rtab}[2]{\par\begin{center}\refstepcounter{table}\begin{minipage}{\linewidth}\centering
  \parbox{\linewidth}{\footnotesize\textbf{Table~\thetable.} #1}\par\smallskip#2\end{minipage}\end{center}}
% each reviewer and the editor letter start a new page; comments follow one another; every block gets a contents entry
\newcommand{\currev}{}
\newcommand{\toeditor}{\FloatBarrier\clearpage\phantomsection\addcontentsline{toc}{section}{Letter to the editor}%
  \section*{Letter to the editor}}
\newcommand{\reviewer}[1]{\FloatBarrier\clearpage\renewcommand{\currev}{#1}\phantomsection
  \addcontentsline{toc}{section}{Reviewer #1}\section*{Reviewer #1}}
\newcommand{\revgeneral}[1]{\FloatBarrier\phantomsection\addcontentsline{toc}{subsection}{#1}%
  \subsection*{#1}}
% a reviewer's general remarks (comment x.0) follow the reviewer heading directly, without a heading or contents entry
\newcommand{\revcomment}[2]{\FloatBarrier\phantomsection\revzero#1..\relax{\label{sec:#1}}%
  {\addcontentsline{toc}{subsection}{#1: #2}\subsection*{#1: #2}\label{sec:#1}}}
\long\def\revzero#1.#2.#3\relax#4#5{\ifnum#2=0 #4\else#5\fi}
\newcommand{\revgeneralpage}[1]{\revgeneral{#1}}
% captions of the figures that later comments link to (each figure is shown once)
\newcommand{\capbaselines}{\textbf{Text-to-brain baselines.} Text-to-brain maps from NeuroVLM and other methods, scored against each PubMed test paper's reconstructed map ($n = 2{,}971$). Panels show (a)~Dice over each map's top 30\% of voxels and (b)~voxelwise Pearson $r$, with means and 95\% bootstrap intervals. Neurosynth, NeuroQuery, and Text2Brain were used as published, and their corpora each include nearly half of the test papers.}
\newcommand{\capdataflow}{\textbf{Data flow of the final model.} Numbered boxes are the trained components, in training order. Gray arrows carry training data or weights, and orange dashed arrows evaluation-only sets. The splits sum to 29,868 because 958 papers with more than 100 peaks or an empty map are in no split. Table~\ref{tab:stages} gives each stage's counts.}
\newcommand{\capbounds}{\textbf{NeuroVault retrieval.} NeuroVLM recall from (a)~text to brain and (b)~brain to text, with each paper's best map (scored highest for its text), worst map (lowest), or every map as target. Chance (dashed) gives each paper's text to another paper at random, applies the same rule, and averages 20 such recall curves. Picking the best of several maps lifts chance above the $k/n$ diagonal even for an unrelated text (AUC 0.74 and 0.70). Picking the worst lowers it (0.31 and 0.32). The every-map rule picks nothing, so chance is the diagonal in (b) but above it in (a), where any of a paper's maps is a hit (0.51 and 0.76).}
\begin{document}
\thispagestyle{empty}
\begin{center}
  {\LARGE\bfseries NeuroVLM: A generative vision-language framework\\[0.15em] for human neuroimaging\par}
  \vspace{0.5em}{\Large Response to reviewers\par}
  \vspace{0.6em}{\large Ryan Hammonds, Jerjes Aguirre-Chavez, Borngreat Omoma-Edosa,\\ Aarya Patel, Bradley Voytek\par}
\end{center}
\hypertarget{contents}{}
\begingroup\small\setlength{\parskip}{0pt}\renewcommand{\contentsname}{Contents}\tableofcontents\endgroup

\toeditor
We thank the editor and reviewers for their thoughtful comments. The revisions and response address every comment. The main changes are:

\begin{itemize}
  \item Independent voxel-space evaluation of text-to-brain maps against simple baselines and existing, published models (\cmt{1.1}).
  \item Retrieval measured by recall@k, mean reciprocal rank, and median rank, also against hard negatives (\cmt{1.1}).
  \item A comparison of generated and retrieved captions (\cmt{3.3}).
  \item A data-flow diagram of the training stages (\cmt{1.2}).
  \item NeuroVault results reported as an upper bound, a lower bound, and a multi-target evaluation, with map types handled separately (\cmt{1.3}).
  \item A leakage audit beyond DOI matching (\cmt{2.4}).
  \item Out-of-distribution and failure-case analyses (\cmt{1.4} and \cmt{2.12}).
  \item Corrected numbers, names, and claims throughout.
\end{itemize}

We also updated the model and its evaluation:

\begin{itemize}
  \item The text-to-brain head is trained with a new loss, latent error plus decoded-map correlation, which raises its Dice (\cmt{1.1} and \cmt{3.5}).
  \item Generated maps are calibrated to the scale of the training maps, with one factor per model fitted on the validation set.
  \item An MLP projector replaces the Qformer in both brain-to-text stages (\cmt{2.1}).
  \item Network and region titles are steered toward canonical names. The language model generates every title, and tokens that cannot continue a canonical name are masked (\cmt{2.12}).
  \item Every evaluation leaves out the duplicated papers found by the leakage audit (\cmt{2.4}).
  \item All result figures were redrawn from the final model at the journal's text width with larger text.
\end{itemize}
% calibration: one global factor per model fitted on the validation set, no per-map fallback (z64/calibration.md; RESPONSE.md hidden note)

Comments are numbered in order within each reviewer, including minor comments and code availability, so \cmt{1.5} is the first minor comment of Reviewer 1. Each comment is quoted in a shaded box and followed by our response and the changes to the manuscript. Comment numbers are links. Figures and tables in this letter are numbered R1, R2, and so on.

Line and page numbers in the changes to the manuscript refer to the revised manuscript with changes marked, which shows additions in blue, deletions in red, and new or changed figures in a blue frame.

\reviewer{1}
\revcomment{1.0}{General remarks}\begin{revquote}
This is an ambitious and timely manuscript introducing \textbf{NeuroVLM}, a unified framework for contrastive text–brain retrieval, text-to-brain generation, and brain-to-text generation. The integration of these capabilities within a modular architecture is appealing, and the use of a relatively large corpus of 30,826 publication-level text–neuroimage pairs is a notable strength. The public release of the code, pretrained models, and datasets should also be valuable to the community. The results provide encouraging evidence that the model learns meaningful associations between scientific text and coarse spatial brain representations. However, I think several central aspects of the evaluation and data construction require further clarification or analysis before the stronger claims regarding map generation and neuroimage interpretation are fully supported.
\end{revquote}\response We added the requested analyses and clarifications (\cmt{1.1} to \cmt{1.10}) and narrowed the claims about map generation and neuroimage interpretation to what the results support (\cmt{1.4}).

\revcomment{1.1}{More independent evaluation of the generative outputs}\begin{revquote}
\textbf{1. More independent evaluation of the generative outputs is needed}
\revqpar The current generative evaluation is performed partly within NeuroVLM’s own learned representation space. Generated maps and text are embedded using the model’s contrastive projection heads and then ranked against candidate items in that same space. This is a useful measure of internal consistency, but it does not independently establish that a generated map is anatomically accurate or that a generated description faithfully reflects the source image.
\revqpar The matched-versus-null results are statistically strong, although some absolute differences are relatively modest—for example, Dice is 0.47 versus 0.40, BERTScore is 0.56 versus 0.52, and semantic cosine is 0.47 versus 0.40. These findings support above-null alignment, but additional validation would help establish neuroscientific specificity.
\revqpar I encourage the authors to add some combination of:
\begin{itemize}
\item Evaluation using independent image and text representations not used during NeuroVLM training.
\item Blinded expert ratings of anatomical accuracy, semantic accuracy, specificity, and unsupported claims.
\item Hard negatives drawn from the same cognitive domain or from studies with similar spatial patterns.
\item Comparison between generated outputs and their nearest training examples.
\item Strong simple baselines, including nearest-neighbor retrieval, mean or domain-specific maps, linear regression from text embeddings to spatial coefficients, and atlas-based labeling.
\end{itemize}
\revqpar It would also be useful to report conventional retrieval measures such as recall@1, recall@5, recall@10, mean reciprocal rank, and median rank, in addition to the area under the full recall@k curve.
\end{revquote}\response We agree and added each evaluation below. Fig.~\ref{fig:baselines} and Table~\ref{tab:t2b} score text-to-brain maps in voxel space against each test paper's reconstructed map. NeuroVLM gives the closest maps (Dice 0.53 against at most 0.49). Its linear-head version is closest behind, then Text2Brain, the mean and domain maps, atlas labeling, Neurosynth, NeuroQuery, and the nearest training paper. Neurosynth, NeuroQuery, and Text2Brain were used as published, and their corpora each hold nearly half of the test papers, which can only favor them. NeuroVLM also leads on the 1,157 test papers in neither corpus.

\textbf{Independent evaluation.} Both directions are scored outside NeuroVLM's representation space. Generated maps are compared in voxel space with the reconstructed map and with the raw coordinate map, which no component was trained on, and NeuroVLM leads on both. Generated captions are compared with the paper's text by MiniLM cosine and BERTScore, whose encoders were not used in training (Reviewer 3, \cmt{3.3}).

\textbf{Modest differences.} The gaps are modest because generated captions describe the map rather than the study, and general-domain scorers compress the range. Within NeuroVLM's latent space, learned from training papers only, matched and shuffled captions separate far more (contrastive cosine 0.28 against 0.01). That space is part of the model, so we report it beside the independent scores, not in their place. Neuroscientific specificity is tested by the hard negatives and the voxel-space baselines.

\textbf{BERTScore.} BERTScore mostly rewards the summary style of the reference, not its content. Random training texts in that style score higher than matched captions (Table~\ref{tab:genret}), and the matched-versus-shuffled gap is only 0.553 against 0.548 when the null is computed from the same captions. The Supplementary Information reports this null.

\textbf{Expert ratings.} We did not collect blinded expert ratings, which are slow and hard to scale, and used existing expert labels and automated checks instead. Anatomical and semantic accuracy are scored by checking region titles against anatomical atlases and network titles against expert atlas labels (Reviewer 3, \cmt{3.6}). Specificity is scored with hard negatives. Unsupported claims are counted as study details that a map cannot support, such as a sample size or patient group (Reviewer 3, \cmt{3.3}).

\textbf{Hard negatives.} Table~\ref{tab:retrieval} scores retrieval among all test papers and among hard negatives, the test papers from the query's own cognitive domain. Hard negatives are only slightly harder than the same number of random papers (recall@10 0.24 against 0.28), so retrieval does not rest on the domain alone.

\textbf{Nearest training examples.} Generated maps match each test paper's map better than the nearest training paper's map does (Table~\ref{tab:t2b}). Generated captions are as close to the paper's text as the nearest training caption by MiniLM cosine, but 57\% of the retrieved captions report details of another study, against 2\% of the generated ones (Reviewer 3, \cmt{3.3}).

\textbf{Simple baselines.} Fig.~\ref{fig:baselines} and Table~\ref{tab:t2b} include every suggested baseline. The linear regression is NeuroVLM with a linear head, fit by ridge regression from the SPECTER2 embedding to the autoencoder latent, which holds the spatial coefficients of each map. The ridge penalty was chosen on the validation set. The final head adds one hidden layer and is trained with the new loss, latent error plus decoded-map correlation, in place of the submission's voxelwise cross-entropy (Table~\ref{tab:ablation}).

\textbf{Retrieval measures.} Table~\ref{tab:retrieval} reports recall@1, recall@5, recall@10, mean reciprocal rank, and median rank in both directions, beside the area under the full recall curve in \msfig{4}. Table~\ref{tab:retrieval} is new Supplementary \mstab{S8}, and Fig.~\ref{fig:baselines} is new Supplementary \msfig{S6}.
\rfig{fig:baselines}{figures/1.1_text_to_brain_baselines.png}{\capbaselines}
% the paper's Fig. 4 (../main/fig4_generative_metrics.py; paired differences: fig4_generative_metrics_text.json)
\rfig{fig:fig4}{figures/fig4_generative_metrics.png}{\textbf{New \msfig{4}.} Text-to-brain (a) and brain-to-text (b) generation by the final model on the PubMed test set ($n = 2{,}971$). Left panels show recall when each generated output retrieves its true item, with $k/n$ the fraction of all candidates. Boxes compare matched (Match, blue) and shuffled pairs (Null, gray) for each score. Stars mark $p < 0.05$, $0.01$, and $0.001$.}
% source: core/a1_text_to_brain_baselines/outputs/results_table_full.csv (PubMed test set, n = 2,971); Text2Brain row: same file (released weights, two fixes for current libraries)
\rtab{\textbf{Text-to-brain maps on the PubMed test set.} Mean Dice and Pearson $r$ between each method's map and the test paper's reconstructed map ($n = 2{,}971$), as in Fig.~\ref{fig:baselines}. The nearest training paper is the one whose text is closest to the query.\label{tab:t2b}}{\small\begin{tabular}{lcc}\toprule
\textbf{Method} & \textbf{Dice} & \textbf{Pearson $r$} \\\midrule
NeuroVLM (non-linear head) & 0.531 & 0.322 \\
NeuroVLM (linear head) & 0.526 & 0.311 \\\midrule
\multicolumn{3}{l}{\textbf{Simple baselines}} \\
\quad Domain mean map & 0.448 & 0.245 \\
\quad Atlas-based labeling & 0.431 & 0.175 \\
\quad Global mean map & 0.420 & 0.251 \\
\quad Nearest training paper & 0.395 & 0.114 \\\midrule
\multicolumn{3}{l}{\textbf{Published models}} \\
\quad Neurosynth & 0.418 & 0.175 \\
\quad NeuroQuery & 0.411 & 0.179 \\
\quad Text2Brain & 0.485 & 0.293 \\\midrule
Shuffled null & 0.408 & 0.122 \\
\bottomrule\end{tabular}}
% source: core/a1_retrieval_metrics_table/outputs/retrieval_metrics_table.tex (final contrastive heads; matches validation/final/report.md)
\rtab{\textbf{Contrastive retrieval on the PubMed test set.} Recall@k, mean reciprocal rank (MRR), and median rank of the final model. Hard negatives are test papers from the query's own cognitive domain, and random negatives match their number.\label{tab:retrieval}}{\small\begin{tabular}{lccccc}\toprule
 & Recall@1 & Recall@5 & Recall@10 & MRR & Median rank \\\midrule
\multicolumn{6}{l}{\textbf{All candidates} (2,971 test papers)} \\
\quad Text $\rightarrow$ brain & 0.018 & 0.053 & 0.088 & 0.045 & 230 \\
\quad Brain $\rightarrow$ text & 0.013 & 0.052 & 0.088 & 0.041 & 217 \\
\quad \textit{Chance} & $<$0.001 & 0.002 & 0.003 & 0.003 & 1,486 \\\midrule
\multicolumn{6}{l}{\textbf{Hard negatives} (papers of the query's own domain, mean 585)} \\
\quad Text $\rightarrow$ brain & 0.054 & 0.160 & 0.242 & 0.118 & 44 \\
\quad Brain $\rightarrow$ text & 0.059 & 0.176 & 0.263 & 0.126 & 40 \\
\quad \textit{Chance} & 0.003 & 0.017 & 0.034 & 0.019 & 250 \\\midrule
\multicolumn{6}{l}{\textbf{Random negatives} (random papers, as many as the hard negatives)} \\
\quad Text $\rightarrow$ brain & 0.069 & 0.196 & 0.283 & 0.142 & 35 \\
\quad Brain $\rightarrow$ text & 0.071 & 0.201 & 0.289 & 0.144 & 35 \\
\quad \textit{Chance} & 0.003 & 0.017 & 0.034 & 0.019 & 250 \\
\bottomrule\end{tabular}}
\revchanges{Methods, Text-to-Brain, \msl{364}{385} (final head), and Performance Metrics, \msl{398}{404} and \msl{415}{427}. Results, Generative Models: PubMed, \msl{596}{613}, and \msfig{4} caption (\msp{15}). Supplementary Information, Projection Head, \msl{1383}{1396} (training of the final head), and Generative Performance, \msl{1506}{1508} (BERTScore null), with \mstab{S4} (\msp{34}). New Supplementary Section Text-to-Brain Baselines and Retrieval, \msl{1935}{1953}, with \msfig{S6} (\msp{44}) and \mstab{S8} (\msp{45}).}

\revcomment{1.2}{Publication-level supervision and dataset separation}\begin{revquote}
\textbf{2. The publication-level supervision and separation of datasets across training stages need clarification}
\revqpar The primary training pairs associate publication titles and abstracts with maps reconstructed from reported coordinate tables. Because a single neuroimaging publication may include multiple experiments, contrasts, populations, and positive or negative effects, the correspondence between an abstract and one aggregated map may be substantially many-to-many.
\revqpar Please clarify:
\begin{enumerate}
\item Whether coordinates were combined across all experiments and contrasts within each publication.
\item How positive and negative effects were handled.
\item How multiple coordinate tables associated with one DOI were treated.
\item The final numbers of publications, experiments, contrasts, and coordinate tables.
\item Whether a contrast-level evaluation can be provided for a subset of the data.
\end{enumerate}
\revqpar The complete train–test separation across model components also needs to be documented more clearly. In particular, Qwen3-0.6B was adapted using 1.2 million neuroscience articles. Please confirm that held-out PubMed and NeuroVault publications were excluded from this language-model corpus. Similarly, the manuscript should explain how test concepts and documents were excluded when constructing the synthetic fine-tuning data and canonical prototype bank. The synthetic stage uses generated maps and descriptions of common networks, regions, and cognitive concepts, creating another possible route through which test information could enter the system.
\revqpar A data-flow diagram showing the datasets used for autoencoder training, contrastive training, language-model adaptation, Qformer training, synthetic fine-tuning, validation, and final testing would substantially improve transparency.
\end{revquote}\response Each map pools every coordinate of one paper. We answer the five points below, confirm the separation of datasets, and added the requested data-flow diagram (Fig.~\ref{fig:dataflow}) and the final counts.

\textbf{1. Experiments and contrasts.} Yes. Each paper's map combines the coordinates of all its experiments and contrasts, because the released coordinates carry no experiment or contrast labels. Each map therefore matches the paper-level text it is paired with.

\textbf{2. Positive and negative effects.} Positive and negative effects are pooled, because the released coordinates carry no sign. Each map shows only where effects lie, not their direction. The Methods state this.

\textbf{3. Multiple coordinate tables.} All coordinate tables of a paper are pooled into one map. Among open-access papers with a coordinate table, 42\% report two or more.

\textbf{4. Final counts.} The released data hold 30,826 publications. Experiments and contrasts cannot be counted, because the released coordinates carry no labels for them. Coordinate tables can be counted only in the open-access full texts, which exist for fewer than half the papers and hold 20,541 tables. The Methods give the splits, and Fig.~\ref{fig:dataflow} the data of each training stage.

\textbf{5. Contrast-level evaluation.} Yes, on NeuroVault, whose maps are released per contrast. We added a per-map evaluation in which each of a paper's maps is scored as its own target against the paper's text (\cmt{1.3}). The model beats chance in both directions (AUC 0.62 and 0.63 against about 0.50).

\textbf{Separation of datasets.} Fig.~\ref{fig:dataflow} is the requested data-flow diagram. No validation or test paper trains the autoencoder, the projection heads, or the projectors. The validation, test, and NeuroVault evaluation papers are excluded from the language model's corpus by its training script. As a memorization check, the longest phrase a generated caption shares with its paper's abstract is two words at the median.

\textbf{Synthetic data and prototype bank.} The synthetic maps are built from training papers only (Reviewer 2, \cmt{2.4}). The prototype bank holds short network, region, and function texts and no test document. Its network entries share names with the evaluation atlas networks, but no atlas map trains any component.
\rfig{fig:dataflow}{figures/1.2_dataflow.png}{\capdataflow}
\revchanges{Methods, Datasets, \msl{137}{144} (pooling per paper and splits), and Training, \msl{169}{181} (synthetic maps from training papers only, exclusions, and the evaluation papers removed from the language-model corpus). Supplementary Information, Coordinate Tables, \msl{1090}{1093}. Discussion, \msl{927}{934}. New Supplementary Section Data Flow and Training Details, \msl{1791}{1799}, with \msfig{S5} (\msp{40}) and \mstab{S6} (\msp{41}).}

\revcomment{1.3}{NeuroVault evaluation}\begin{revquote}
\textbf{3. The NeuroVault evaluation would benefit from target-independent selection and map-type-specific preprocessing}
\revqpar The supplementary methods state that when several NeuroVault images were associated with the same title and abstract, the contrastive model selected the most similar image. The manuscript appropriately acknowledges that this inflates performance and provides an upper bound.
\revqpar I suggest retaining this as an exploratory upper-bound analysis while adding a target-independent evaluation. For example, the authors could use contrast-specific NeuroVault metadata, evaluate all eligible images, apply a predefined image-selection rule, or use a multiple-instance retrieval metric.
\revqpar The preprocessing of the NeuroVault images also requires reconsideration. The dataset contains t-, z-, F-, p-, and beta maps, as well as ROI masks and parcellations, but all maps are processed with a 99th-percentile threshold, cluster filtering, and binarization. These map types have substantially different meanings. In particular, the largest values in a conventional p-value map generally indicate weaker evidence, signed t-, z-, and beta maps may contain meaningful negative effects, and parcellations represent categorical labels rather than statistical intensity.
\revqpar Please clarify whether map-specific transformations were applied before thresholding. If not, I recommend separate preprocessing and reporting for signed statistical maps, p-maps, masks, and parcellations. This would make the NeuroVault generalization results easier to interpret and ensure that spatial information is treated appropriately for each map type.
\end{revquote}\response We kept the model-selected map as an upper bound and added two target-independent rules and a lower bound (Fig.~\ref{fig:bounds}). Multi-target retrieval counts any of a paper's maps as a hit. The per-map score makes each map its own target. The lower bound keeps each paper's worst-matching map. The model beats chance under every rule, and per-map retrieval reaches an area under the recall curve (AUC) of 0.62 from text to brain and 0.63 from brain to text.

\textbf{Map types.} No map-specific transformation was applied before thresholding in the submission, which thresholded every map alike after removing masks, parcellations, anatomical maps, and variance maps. P-maps are inverted to 1 $-$ p, so that high values mark strong evidence. Signed maps keep their positive tail, as before, because using both tails, or adding negative tails as extra maps, reduces evaluation performance. Table~\ref{tab:maptypes} reports each map type separately, and \msfig{5}b uses the same handling.

\rfig{fig:bounds}{figures/1.3_neurovault_bounds.png}{\capbounds}
\rtab{\textbf{NeuroVault results by map type.} AUC of the final model when every map is a target, scored per map. From text to brain, a paper's text ranks each of its maps among the maps of other papers. From brain to text, each map ranks the 298 texts, as in Fig.~\ref{fig:bounds}b. Chance is this setting's null, texts shuffled across papers (0.50 and 0.51). Rows beginning with Other are maps that their uploaders labeled Other, which we typed from their file names, descriptions, and values. Masks, parcellations, anatomical and variance maps, mean images, and untyped maps were excluded.\label{tab:maptypes}}{\small\begin{tabular}{L{3.6cm}rrlcc}\toprule
 & & & & \multicolumn{2}{c}{\textbf{AUC}} \\\cmidrule(lr){5-6}
\textbf{Map type} & \textbf{Maps} & \textbf{Papers} & \textbf{Handling} & \textbf{text $\rightarrow$ brain} & \textbf{brain $\rightarrow$ text} \\\midrule
T & 1,761 & 189 & positive tail & 0.64 & 0.64 \\
Z & 811 & 75 & positive tail & 0.57 & 0.58 \\
beta, univariate & 43 & 5 & positive tail & 0.61 & 0.77 \\
F & 40 & 17 & positive tail & 0.64 & 0.55 \\
beta, multivariate & 14 & 5 & positive tail & 0.84 & 0.83 \\
p & 3 & 2 & inverted (1 $-$ p) & 0.64 & 0.60 \\
1 $-$ p & 2 & 1 & used as is & 0.54 & 0.44 \\
Other, t, z, or F & 204 & 23 & positive tail & 0.63 & 0.70 \\
Other, effect estimate & 103 & 15 & positive tail & 0.69 & 0.62 \\
Other, component or gradient & 93 & 8 & positive tail & 0.49 & 0.51 \\
\bottomrule\end{tabular}}
\revchanges{\msfig{5}b caption (\msp{17}). Results, Contrastive Model: NeuroVault, \msl{803}{816}. Supplementary Information, NeuroVault preprocessing, \msl{1229}{1246} (map types and exclusions) and \msl{1257}{1263} (bounds). New Supplementary Section NeuroVault Bounds and Map Types, \msl{2046}{2055}, with \mstab{S9} (\msp{46}).}

\revcomment{1.4}{Scope of the generative claims}\begin{revquote}
\textbf{4. The scope of the generative claims should more closely reflect the available training targets}
\revqpar The authors appropriately acknowledge that the training images are proxy maps reconstructed from reported peak coordinates. Such maps omit subthreshold structure, cluster shape, effect size, uncertainty, contrast definitions, and other study-specific information. The additional smoothing, DiFuMo reconstruction, resampling, and binarization further emphasize broad spatial support rather than complete statistical information.
\revqpar Given these limitations, the generated outputs may be more accurately described as \textbf{coarse spatial-support predictions} or \textbf{spatial priors} rather than complete statistical maps. This framing would still support useful applications in literature navigation, network labeling, and exploratory hypothesis generation, while avoiding the implication that effect size, sign, uncertainty, or a study-specific statistical contrast has been reconstructed.
\revqpar Similarly, the brain-to-text outputs appear designed to provide concise, high-level network or functional descriptions rather than detailed interpretations of individual experiments. The manuscript should distinguish broad network labeling from study-specific interpretation and mechanistic inference.
\revqpar The authors already note that the model can generate plausible maps for essentially any text input and should not be viewed as an oracle. A small quantitative analysis of failure cases, out-of-distribution inputs, and abstention or confidence calibration would strengthen this important caution.
\end{revquote}\response We agree and adopted this framing. The Abstract, Results, and Discussion describe text-to-brain outputs as coarse spatial priors that carry no effect size, sign, or uncertainty, and brain-to-text outputs as broad network-level labels, not interpretations of a specific study or mechanism.

\textbf{Failure cases.} In 16\% of the 2,971 PubMed test-set captions, the similarity to the paper's text is negative. Many of these failures start in the data, where the paper's own text and map are already among the least similar pairs of the test set. The rest are model misses. Our response to Reviewer 2 analyzes them (\cmt{2.12}).

\textbf{Out-of-distribution inputs.} Warning scores flag inputs far from the training data (Fig.~\ref{fig:ood}). A map is flagged when it correlates poorly with its autoencoder reconstruction or when its latent lies far from the training latents, and a text when it lies far from every training abstract. With thresholds set on the validation split, the scores flag about 5\% of PubMed test-set maps and abstracts, every shuffled, white-noise, and rescaled map, every off-topic text, and 61 to 80\% of rotated or spun maps.

\textbf{Abstention and confidence.} Flagged inputs can be withheld, but the warning scores are not calibrated probabilities. Within the training distribution, the peak of a calibrated text-to-brain map serves as a confidence score, because it tracks the map's retrieval rank (Spearman 0.63). The peak does not detect off-topic text, whose maps mostly peak within the range of the test abstracts.
% produced by core/a7_ood_confidence/make_figure.py (outputs/flagged.csv; thresholds on the validation split, 2,961 maps / abstracts: map r(map, reconstruction) < 0.46 (2.5th pct) or latent Mahalanobis > 27.0 (97.5th pct), text SPECTER2 cosine to nearest training abstract < 0.79 (5th pct). Flagged: validation maps 5.0 %, test maps 4.5 %, rot90 / rotation / spin 61 / 70 / 80 %, shuffle, noise, x100, x0.01 100 %; test abstracts 5.2 %, WikiText and recipes 100 %; AUROC 0.89 to 1.00. The previous plain cosine missed white noise (0 %): a dense non-negative map and its flat reconstruction share a constant component, cosine 0.63 against a 0.55 threshold, r 0.00 (outputs/diagnosis.csv). Single scores at the validation 95th percentile: outputs/detection.csv. 16 % = test-set captions with a negative similarity to the paper's text (2.12, a7_failure_cases_percentiles))
\rfig{fig:ood}{figures/1.4_ood.png}{\textbf{Out-of-distribution warnings.} Scores of PubMed test-set maps and abstracts (blue) and of altered maps and off-topic texts (gray). Panels show (a)~each input's similarity to the training data, (b)~the Mahalanobis distance of a map's latent from the training latents, and (c)~the share of inputs flagged. Dashed lines mark thresholds set on the validation split.}
\revchanges{Abstract, \msl{28}{39}. Introduction, \msl{112}{117}. Methods, Text Decoder, \msl{320}{327} (out-of-distribution warnings), and Text-to-Brain, \msl{374}{385}. Results, Generative Models, \msl{433}{435} (spatial priors, network-level labels). Discussion, \msl{909}{910}, and limitations, \msl{948}{962}. New Supplementary Section Out-of-Distribution Inputs, \msl{2091}{2104}, with \msfig{S7} (\msp{47}).}

\revcomment{1.5}{Inconsistent dataset and metric values}\begin{revquote}
1. \textbf{Please reconcile inconsistent dataset and metric values.} The main methods list 3,000 NeuroVault maps and 100 network images, whereas the supplement reports 4,553 NeuroVault maps and the results describe 300 network images. The autoencoder AUROC is also reported as 0.94 in the Figure 5 caption and 0.975 in the main text.
\end{revquote}\response We recomputed every count from the released data and corrected the text and captions.

\textbf{NeuroVault maps.} The Methods' 3,000 was a rounded count, and the supplement's 4,553 maps predate image quality control and the removal of papers already in PubMed. After these steps, the leakage audit, and the removal of mean, variance, and untyped maps, retrieval uses 3,074 maps.

\textbf{Network images.} We corrected both counts to the 185 released network maps. They are the atlas maps and the ICA components.

\textbf{Autoencoder AUROC.} Both values were correct but referred to different datasets, the caption's 0.94 to NeuroVault and the main text's 0.975 to the PubMed test set. With the final model, the \msfig{5} caption and the main text give 0.939 on NeuroVault and 0.974 on the PubMed test set, and each names its dataset.
\revchanges{Methods, Datasets, \msl{132}{144}. Results, Contrastive Model: Network Labeling, \msl{785}{797}, and Autoencoder, \msl{876}{880}. Figs.~\msfign{4} and \msfign{5} captions (\mspp{15}{17}). Supplementary Information, NeuroVault preprocessing, \msl{1229}{1246} and \msl{1262}{1264}.}

\revcomment{1.6}{SPECTER and SPECTER2}\begin{revquote}
2. \textbf{Please use SPECTER and SPECTER2 consistently.} Section 2.6 describes SPECTER, whereas the text-to-brain methods and supplement describe SPECTER2.
\end{revquote}\response Every text-encoding step uses SPECTER2, the allenai/specter2\_aug2023refresh model with its ad hoc query adapter. The original SPECTER is never used. We corrected every mention to SPECTER2.
\revchanges{Introduction, \msl{107}{111}. Methods, Text Encoder, \msl{220}{227}, and Text-to-Brain, \msl{364}{373}. Supplementary Information, Text preprocessing, \msl{1131}{1134}. \msfig{S1} caption (\msp{24}).}

\revcomment{1.7}{Two methodological analogies}\begin{revquote}
3. \textbf{Two methodological analogies should be revised.} Transformer language models still depend on a fixed tokenizer vocabulary, despite their ability to represent context and novel compositions. In addition, a separately trained projection head is not technically equivalent to LoRA, although both approaches can preserve the weights of a larger pretrained encoder.
\end{revquote}\response We agree with both points. The tokenizer sentence states that contextual encoders also depend on a fixed vocabulary and that their advantage is composing words in context. The LoRA analogy is deleted, and the head is named the text-to-brain head. Two analyses for Reviewer 3 test the composition claim. SPECTER2 beats a context-free word2vec encoder (\cmt{3.4}), and swapping one word moves the generated map as expected, although negation has no effect (\cmt{3.5}).
\revchanges{Introduction, \msl{81}{84} (tokenizer sentence). Methods, Text-to-Brain, \msl{364}{385} (LoRA sentence deleted, projection head named).}

\revcomment{1.8}{Edge-device claim}\begin{revquote}
4. \textbf{The edge-device claim would benefit from empirical support.} Please report total memory requirements, inference latency, hardware, numerical precision, and any quantization used. The relatively small text encoder is attractive, but the complete brain-to-text pipeline also includes the Qformer and Qwen language model.
\end{revquote}\response Table~\ref{tab:edge} reports memory, latency, hardware, and numerical precision for both full pipelines, projector and language model included, on a consumer GPU and a desktop CPU. On the CPU, dynamic int8 quantization runs brain-to-text fastest but saves no memory over bf16. We tested no phone or embedded device, so the Methods claim a consumer CPU instead of edge devices. The timings use the released weights, whose projector and text-to-brain head have the same size as the final ones.
\rtab{\textbf{Cost of one query.} Median latency at batch 1 and peak memory of weights and activations, not file size. Brain-to-text decodes 128 tokens greedily. The GPU is an RTX 5090 and the CPU an 8-thread Ryzen 7 9800X3D. Dynamic int8 quantization leaves the token embedding, the tied output layer, and SPECTER2's adapter-wrapped layers in bf16, so it saves no memory.\label{tab:edge}}{\small\begin{tabular}{lrcccc}\toprule
\textbf{Pipeline, one query} & \textbf{Parameters} & \textbf{GPU, bf16} & \textbf{CPU, fp32} & \textbf{CPU, bf16} & \textbf{CPU, int8} \\\midrule
text-to-brain & 171 M & 7 ms, 0.4 GB & 135 ms, 0.6 GB & 53 ms, 0.4 GB & 120 ms, 0.5 GB \\
brain-to-text, 128 tokens & 685 M & 3.2 s, 1.3 GB & 13.1 s, 2.7 GB & 6.3 s, 1.4 GB & 3.5 s, 1.4 GB \\
\bottomrule\end{tabular}}
\revchanges{Methods, Text Encoder, \msl{221}{230} (edge-device sentence replaced by CPU costs). Discussion, \msl{988}{989}.}

\revcomment{1.9}{Figure text size, full descriptions, affiliations}\begin{revquote}
5. Some text in Figures 2, 3, and S4 is difficult to read at normal manuscript scale. Providing the complete generated descriptions in a supplementary table would be helpful. Please also correct the remaining “Organization” placeholders in the affiliation block.
\end{revquote}\response We redrew Figs.~\msfign{2}, \msfign{3}, and \msfign{S4} at the journal's text width with larger text (Figs.~\ref{fig:fig2} to \ref{fig:figS4}). New Supplementary \mstab{S5} gives the complete generated texts of \msfig{2}, and Figs.~\msfign{3} and \msfign{S4} print each generated text in full beside its map. The affiliation placeholders are corrected.
\rfigparts{fig:fig2}{figures/1.9_fig2_network_examples}{\textbf{New \msfig{2}.} Eight atlas networks, each with its atlas map (True), the map generated from its description, and the network, anatomy, and cognition texts generated from the atlas map. Auditory and default mode show the complete texts, and the other six their titles. Maps show their top 5\% of voxels in sagittal, coronal, and axial views. Supplementary \mstab{S5} gives every text.}
\rfig{fig:fig3}{figures/1.9_fig3_newgen.png}{\textbf{New \msfig{3}.} Six NeuroVault papers, each with its title and the map generated from its text, then its true map and the text generated from that map. The distributions of both shared-space cosines are drawn on the left, and the rows are drawn from their high, medium, and low thirds by a fixed rule (Reviewer 3, \cmt{3.9}). Low rows show a scattered true map (top) and a model miss on a clean map (bottom). Maps show their top 5\% of voxels.}
\rfig{fig:figS4}{figures/1.9_figS4_clip_ica.png}{\textbf{New \msfig{S4}.} Labeling of resting-state components from UK Biobank (top) and the Human Connectome Project (bottom). Bars give the five most similar labels from the contrastive models. The generated column is the {[}NETWORK{]} caption of the stage 2 MLP projector, its title in bold.}
\revchanges{Affiliations, \msl{16}{17}. Figs.~\msfign{2}, \msfign{3}, and \msfign{S4}, redrawn (\msp{11}, \msp{12}, and \msp{35}). Results, Generative Models: Networks, \msl{443}{449}. New Supplementary \mstab{S5} with the complete texts of \msfig{2}, \msl{1613}{1616} (\mspp{36}{39}). Figs.~\msfign{4} and \msfign{5} captions (\mspp{15}{17}).}

\revcomment{1.10}{Overall assessment}\begin{revquote}
NeuroVLM is a creative and potentially useful framework, and the current results provide encouraging evidence that large collections of publication text and coordinate-derived maps can support cross-modal retrieval, broad functional labeling, and coarse generation. The manuscript would be considerably strengthened by independent validation of the generative outputs, clearer separation of datasets across all training stages, a revised NeuroVault evaluation, and more precise framing of the generated spatial representations. These revisions would help establish the appropriate scientific use cases and limitations of the framework.
\end{revquote}\response The four revisions named here are made in our responses to \cmt{1.1} to \cmt{1.4}, in the same order.

\revcomment{1.11}{Code availability}\begin{revquote}
The repository is a useful community resource: it has an Apache-2.0 license, pip installation instructions, pretrained models and data, tutorials, tests, and CI across Python 3.10–3.13 on macOS and Ubuntu. The latest CI run passed 405 tests (one skipped), although reported code coverage was 47\%.
\revqpar The paper's results appear only partially reproducible. The repository includes data, training, and evaluation notebooks, but two preprocessing notebooks contain unresolved merge conflicts and are invalid; some analyses rely on manual steps, machine-specific paths, or external checkpoints; dependencies and Hugging Face artifacts are not fully pinned; and there is no figure-by-figure reproduction guide. Thus, the package appears usable for inference, but exact end-to-end reproduction would require additional fixes and documentation.
\revqpar I reviewed the repository and existing CI results but did not install or run the code locally.
\end{revquote}\response The results reproduce end to end. One command builds the pinned environment, fetches the Hugging Face inputs at pinned and checksummed revisions, retrains every model, redraws the figures, and checks each metric against the paper. No step is manual and no path is machine-specific. Two runs from scratch gave byte-identical weights on the same GPU model and library versions. We resolved the merge conflicts in both notebooks and added a figure-by-figure guide that names the script behind every paper figure.

\textbf{Tests and coverage.} Coverage is [author: new percentage] on the 58 files that train the models and draw the figures. The CNN and other experiments outside the paper are excluded from the measure. Because training takes hours, continuous integration trains every model for one batch on CPU, which exercises the whole pipeline.

\textbf{Remaining limits.} Two inputs remain outside the pinned set, the Automated Anatomical Labeling (AAL) atlas, whose download server fails, and the NiCLIP and NeuroConText outputs compared in \msfig{S3}. [author: one clause on how each is obtained, e.g. "both ship with the repository" or "both are fetched from their authors' releases".]
\revchanges{New Data availability and Code availability sections after the Discussion, \msl{1007}{1039}, replacing the links in Methods, Training, \msl{166}{168}.}

\reviewer{2}
\revcomment{2.0}{General remarks}\begin{revquote}
The manuscript presents NeuroVLM, a unified framework for neuroimage–text contrastive learning, text-to-brain generation, and brain-to-text generation. The topic is meaningful, and the large-scale dataset provides potential value. However, the manuscript is not suitable for publication in its current form and several issues should be addressed.
\end{revquote}\response We addressed each issue (\cmt{2.1} to \cmt{2.12}).

\revcomment{2.1}{Motivation of the Qformer}\begin{revquote}
1.The authors use a Q-Former for brain-to-text generation, but its motivation is unclear. Comparisons with simpler alternatives, such as an MLP projector or cross-attention module, should be provided.
\end{revquote}\response We agree. The Qformer attended over only two input tokens, the map latent and its semantic vector, so its attention had little to select. We compared it with a multilayer perceptron (MLP) projector of the same size, which raises stage 1 retrieval AUC from 0.82 to 0.90 (Fig.~\ref{fig:mlp}). The MLP projector replaces the Qformer in both stages. A cross-attention module would also attend over only two tokens, so we did not test one.
\rfig{fig:mlp}{figures/2.1_mlp_projector.png}{\textbf{Qformer and MLP projector.} Gray bars show the original model with the Qformer and blue bars the final model with the MLP projector, on PubMed test maps. Panels show (a)~retrieval AUC, (b)~caption-to-map cosine, and (c)~BERTScore F1 against the paper's text.}
\revchanges{Methods, Text Decoder, \msl{284}{297}, and Brain-to-Text, \msl{341}{361} (MLP projector). Results, Generative Models: Networks, \msl{443}{449} (projector named). Supplementary Information, Brain-to-Text Projector, \msl{1405}{1425} and \msl{1465}{1474} (rewritten for the MLP projector), with \mstab{S3} (\msp{33}). Figs.~\msfign{S1} and \msfign{S4} captions (\msp{24} and \msp{35}).}

\revcomment{2.2}{Datasets and sample sizes per model}\begin{revquote}
2.The training data for the contrastive and generative models are not clearly described. The authors should specify the datasets and sample sizes used for each model and clarify whether the same image-text pairs were reused across different training stages.
\end{revquote}\response Fig.~\ref{fig:dataflow} and Table~\ref{tab:stages} give the data and sample size of each training stage, and several stages reuse the same training papers. Reviewer 1 raised the same concern (\cmt{1.2}).

\textbf{1. Datasets and sample sizes.} Table~\ref{tab:stages} lists the training and validation size of each component, and Fig.~\ref{fig:dataflow} shows which set feeds each stage.

\textbf{2. Reuse across stages.} Yes. The same training papers and their maps train the autoencoder, the contrastive projection heads, and the stage 1 projector. The anatomically centered head and the stage 2 projector train on short concept texts paired with maps built from training papers only (\cmt{2.6}). No validation or test paper trains any component, and the language model's training script removes them from its corpus (\cmt{2.11}).
% sources: core/a4_counts_table/outputs/stage_splits.md (from scripts/outputs/runs/final/*.summary.json) and
% core/a4_dataflow_diagram/outputs/dataflow_counts.json (same numbers); NeuroQwen3-0.6B split from z53 (01_dapt.ipynb)
\rtab{\textbf{Training data per component.} Training and validation sizes of each trained component of the final model, as in Fig.~\ref{fig:dataflow}. The text-to-brain head is selected on 2,961 validation papers, the 2,987 minus the 26 removed by the leakage audit.\label{tab:stages}}{\small\begin{tabular}{p{4.6cm}p{5.4cm}rrr}\toprule
\textbf{Training stage} & \textbf{Data} & \textbf{Train} & \textbf{Validation} & \textbf{Unit} \\\midrule
Autoencoder & PubMed coordinate maps & 23,894 & 2,987 & papers \\
Contrastive heads & SPECTER2 text and map latents & 23,894 & 2,987 & papers \\
Text-to-brain head & SPECTER2 text and map latents & 23,894 & 2,961 & papers \\
Anatomically centered projection head & synthetic concept text--map pairs from training papers & 6,309 & 64 & pairs \\
Stage 1 projector (MLP) & map latents and LLM summaries & 22,547 & 2,837 & papers \\
Stage 2 projector (MLP) & {[}NETWORK{]}, {[}REGION{]}, {[}FUNCTION{]} texts & 972 & 134 & texts \\
NeuroQwen3-0.6B & PubMed titles and abstracts & 1.22 M & 12,317 & articles \\
\bottomrule\end{tabular}}
\revchanges{Methods, Training, \msl{159}{173}. New Supplementary Section Data Flow and Training Details, \msl{1791}{1799}, with \msfig{S5} (\msp{40}) and \mstab{S6} (\msp{41}), as for Reviewer 1's comment \cmt{1.2}.}

\revcomment{2.3}{NeuroVault selection bias}\begin{revquote}
3.The NeuroVault evaluation has selection bias. When multiple images correspond to the same title and abstract, the contrastive model selects the most similar image. This procedure may substantially inflate the reported performance.
\end{revquote}\response We agree that choosing the most similar map inflates the score. A lower bound keeps each paper's worst-matching map, a multi-target score counts any of the paper's maps as a hit, and a per-map score makes each map its own target (Fig.~\ref{fig:bounds}). Each rule has its own chance curve, from texts shuffled across papers, and the model's AUC stays about 0.1 above chance under every rule. Reviewer 1 raised the same concern (\cmt{1.3}).
\revchanges{\msfig{5}b caption (\msp{17}). Results, Contrastive Model: NeuroVault, \msl{803}{816}. Supplementary Information, NeuroVault preprocessing, \msl{1229}{1246} (map types and exclusions) and \msl{1257}{1263} (bounds). New Supplementary Section NeuroVault Bounds and Map Types, \msl{2046}{2055}, with \mstab{S9} (\msp{46}), as for Reviewer 1's comment \cmt{1.3}.}

\revcomment{2.4}{Data leakage}\begin{revquote}
4.Potential data leakage across PubMed, Neurosynth Compose, and NeuroVault has not been sufficiently excluded. DOI-based filtering alone may not identify duplicate records, preprint and journal versions, or repeated coordinate tables.
\end{revquote}\response We audited leakage beyond DOI filtering with the checks below, which cover duplicate records, preprint and journal versions, and repeated coordinate tables. No DOI or PubMed identifier is shared across splits. The checks flag 16 test and 26 validation papers (Table~\ref{tab:leakage}), and every evaluation leaves them out. Seven NeuroVault papers that duplicate a training paper were also removed.

\textbf{1. Duplicate records and preprint versions.} Each held-out paper was compared with every training and NeuroVault paper by fuzzy title matching and by the similarity of its title and abstract. A preprint and its journal version share a near-identical title and abstract, so this check catches them.

\textbf{2. Repeated coordinate tables.} The coordinate tables from PubMed and Neurosynth Compose were compared across all splits. Identical or near-identical tables mark a paper that repeats another paper's peaks.

\rtab{\textbf{Leakage audit.} Held-out PubMed papers flagged by each check against other papers. Rows overlap, and the last row is their union.\label{tab:leakage}}{\small\begin{tabular}{lcc}\toprule
\textbf{Check against other papers} & \textbf{PubMed test set (2,987)} & \textbf{PubMed validation set (2,987)} \\\midrule
Near-identical title & 4 & 6 \\
Near-identical title and abstract & 3 & 6 \\
Identical coordinate table & 6 & 12 \\
Near-identical coordinate table & 12 & 21 \\
Removed from every evaluation & 16 & 26 \\
\bottomrule\end{tabular}}
\revchanges{Methods, Training, \msl{169}{181} (synthetic maps from training papers only, leakage audit, evaluation exclusions, and the evaluation papers removed from the language-model corpus). Methods, Datasets, \msl{142}{144} (NeuroVault papers that duplicate a training paper removed). Supplementary Information, NeuroVault preprocessing, \msl{1229}{1246}.}

\revcomment{2.5}{Competitive baselines}\begin{revquote}
5.The manuscript lacks competitive baselines for the generative tasks. The proposed methods should be compared with existing approaches and simpler projection or retrieval-based baselines.
\end{revquote}\response NeuroVLM matches or beats existing approaches and simpler projection and retrieval baselines in both directions. Reviewer 1 raised the same concern, and the analyses are in our response to \cmt{1.1}.

\textbf{Text to brain.} Fig.~\ref{fig:baselines} and Table~\ref{tab:t2b} compare NeuroVLM's maps with three published models and four simple baselines. The projection baseline is the same model with a linear head, and the retrieval baseline is the nearest training paper's map. NeuroVLM leads both, as well as Neurosynth, NeuroQuery, the mean and domain maps, and atlas labeling. Text2Brain comes closest among the published models (Dice 0.49 against 0.53) and trails by more on the papers outside its own training corpus.

\textbf{Brain to text.} We know of no published model that writes text from group-level brain maps. The existing decoders NiCLIP and NeuroConText rank a fixed set of texts, and NeuroVLM is slightly ahead of both at that task (\msfig{S3}). The simpler projection baseline, an MLP projector, beats the Qformer and replaces it (\cmt{2.1}). The retrieval baseline, the nearest training caption, is no closer to the paper's text than the generated caption and often carries another study's details (Reviewer 3, \cmt{3.3}).
\revchanges{Methods, Text-to-Brain, \msl{364}{385} (final head), and Performance Metrics, \msl{398}{404} and \msl{415}{427}. Results, Generative Models: PubMed, \msl{596}{613}, and \msfig{4} caption (\msp{15}). New Supplementary Section Text-to-Brain Baselines and Retrieval, \msl{1935}{1953}, with \msfig{S6} (\msp{44}) and \mstab{S8} (\msp{45}), as for Reviewer 1's comment \cmt{1.1}.}

\revcomment{2.6}{Summarized and synthetic text}\begin{revquote}
6.The construction of the summarized and synthetic text data is unclear. The authors should describe the generation procedure and provide representative examples.
\end{revquote}\response One prompt generated the paper summaries and another generated the synthetic and curated texts. Both prompts are given below, and Table~\ref{tab:examples} gives a complete example of each kind of text.

\textbf{Summaries.} Qwen3.5-27B summarized each paper with the first prompt as its system message and the paper's title and abstract, joined by a separator token, as the user message. The prompt asks for a summary of the paper without its study details. Decoding used sampling with temperature 0.3 and top-p 0.9. Re-running the prompt on PubMed test-set papers gives summaries on the same topics as the released ones. The stage 1 projector learns from these summaries.

\textbf{Synthetic texts.} Qwen3.5-27B wrote short entries with the second prompt, from one seed per request, such as ``Insula''. The seeds were region names from standard brain atlases, Cognitive Atlas concepts, and the network names of our network atlases. Decoding used temperature 0.1 and top-p 0.9 with thinking mode off, and each seed was sampled several times to obtain re-wordings of one concept. The anatomically centered head learns from these entries, paired with maps built from training papers (\cmt{2.4}).

\textbf{Curated texts.} The stage 2 projector learns from 1,106 curated texts about networks, regions, and functions. We generated them with Qwen3.5-27B in the same way. A script in the code repository repeats the process, and its texts match the released texts of the same seeds in content (MiniLM cosine 0.88 on average).
% re-creation (z97): scripts/data/08_synthetic_texts.py (seeds: AAL SPM12 + Harvard-Oxford region names, Cognitive Atlas terms,
% `networks` dataset names; examples and cosines in scripts/data/examples/synthetic_texts_examples.md) and
% scripts/data/09_summaries.py (examples in scripts/data/examples/summaries_examples.md)
% the two prompts, verbatim (the two longest lines of the second are wrapped to fit the box; the words are unchanged);
% first box: system prompt and sampling settings from docs/02_data/07_transform_text.ipynb cells 3-4 (z72), model
% Qwen3.5-27B (author, final round; the pubmed_summary_qa dataset card still names Llama-3.3-70B-Instruct)
\begin{center}\begin{lrbox}{\revbox}\begin{minipage}{\dimexpr\linewidth-2\fboxsep-2\fboxrule\relax}\scriptsize
\begin{verbatim}
Model: Qwen/Qwen3.5-27B
Decoding: sampling, temperature 0.3, top-p 0.9, top-k 50, at most 256 new tokens
User message: title[SEP]abstract
System message:

You are writing a short, topic-focused neuroscience mini-review.

INPUT:
You will receive:
    title + " [SEP] " + abstract from an fMRI publication.

GOAL:
Write a single coherent paragraph (4–6 sentences) that:
- Describes the scientific topic implied by the paper (e.g., meditation and the default mode
  network, exercise and working memory), not the paper’s methods.
- Explains the key cognitive processes (e.g., working memory, self-referential thought,
  emotion regulation) and the main brain systems or networks involved.
- May mention typical effects or relationships in general terms (e.g., “has been associated
  with changes in connectivity” or “can influence executive control”), but should stay high-level.

ALLOWED:
- General statements about effects or relationships (e.g., “exercise has been associated
  with improved cognition”, “meditation can alter DMN connectivity”).
- Generic references to “studies” or “research” in the field.
- Qualitative descriptions of brain activity or connectivity (e.g., “increased activation
  in prefrontal regions”, “changes in functional connectivity”).

AVOID:
- Exact numerical details (sample sizes, p-values, coordinates, time points).
- Detailed task or protocol descriptions from the specific study.
- Directly quoting or closely paraphrasing long sentences from the abstract.
- Overly narrow focus on this single experiment; write as if summarizing the topic area.

STYLE:
- Neutral, scientific tone, like a short review or encyclopedia entry.
- One paragraph only. No bullet points or headings.
- Focus on giving a clear sense of “what this topic is about” in the context of the field.
\end{verbatim}
\end{minipage}\end{lrbox}\fbox{\usebox{\revbox}}\end{center}
\begin{center}\begin{lrbox}{\revbox}\begin{minipage}{\dimexpr\linewidth-2\fboxsep-2\fboxrule\relax}\scriptsize
\begin{verbatim}
Write a short educational neuroscience entry for the provided fMRI concept.

Requirements:
- Explain the concept in general neuroscience terms.
- Always include relevant anatomical regions.
- Connect the anatomy to the cognitive or neural function.
- Use cautious, broad language; avoid overclaiming.
- Do not mention studies, patients, diseases, disorders, treatments, age groups, tasks, results,
  methods, or statistics.
- Do not use: critical for, responsible for, causal, proven, this study, activation, altered,
  reduced, increased.
- Do not invent highly specific subregions unless they are commonly associated with the concept.
- If the concept is broad, name the main commonly associated regions or networks.
- If the concept is anatomical, describe its broader network/function associations.
- Write 3-5 complete sentences.

Return valid JSON only:
{
  "title": "1-5 words",
  "description": "3-5 sentences"
}
\end{verbatim}
\end{minipage}\end{lrbox}\fbox{\usebox{\revbox}}\end{center}
\rtab{\textbf{Examples of the training texts.} Each row gives one complete text and the training stage that uses it. Qwen3.5-27B wrote the summaries and the synthetic and curated texts.\label{tab:examples}}{\small\begin{tabular}{p{3.6cm}p{12.3cm}}\toprule
\textbf{Text (used by)} & \textbf{Complete example} \\\midrule
Summary of a training paper (stage 1 projector) & Theory of Mind (ToM) and empathy are crucial processes in social cognition, enabling individuals to attribute mental states and infer emotional experiences in others. These processes involve a network of brain regions, including the medial prefrontal cortex, superior temporal sulcus, and temporal pole. ToM and empathy share overlapping but distinct neuronal networks, with common areas of activation including the medial prefrontal cortex and temporoparietal junction. \\
Synthetic text ``Insula'' (anatomically centered head) & The insula is a deep cortical region situated within the lateral sulcus of the cerebral hemisphere. This area maintains connections with structures such as the anterior cingulate cortex and the amygdala. Processes occurring in this region are often linked to perceptions of internal bodily states and emotional awareness. It contributes to the integration of sensory information with subjective experience across various brain networks. \\
Curated {[}NETWORK{]} text ``Salience Network'' (stage 2 projector) & The salience network commonly includes anterior insula and dorsal anterior cingulate cortex. It is associated with detecting and prioritizing behaviorally relevant internal and external information. Its connections with limbic, autonomic, and frontoparietal systems support flexible shifts between internal monitoring and externally oriented control. \\
Curated {[}FUNCTION{]} text ``Language Comprehension'' (stage 2 projector) & Language comprehension commonly involves superior temporal cortex, middle temporal cortex, inferior frontal cortex, angular gyrus, and anterior temporal cortex. These regions support speech and text meaning, syntactic structure, and semantic integration. Left-hemisphere regions are often prominent, although bilateral temporal systems can contribute to broader language understanding. \\
\bottomrule\end{tabular}}
\revchanges{Methods, Brain-to-Text, \msl{341}{361}. Supplementary Information, Text preprocessing, \msl{1141}{1226} (both prompts, their models, seeds, and settings, and the examples).}

\revcomment{2.7}{Dataset sizes}\begin{revquote}
1.The reported dataset sizes are inconsistent. The Methods section reports 3,000 NeuroVault images and 100 network images, whereas other sections report 4,553 and 300, respectively.
\end{revquote}\response We recomputed every count from the released data and corrected the text and captions. Reviewer 1 raised the same concern (\cmt{1.5}).

\textbf{NeuroVault maps.} The supplement's 4,553 maps predate image quality control and the removal of papers already in PubMed. After these steps, the leakage audit, and the removal of mean, variance, and untyped maps, retrieval uses 3,074 maps.

\textbf{Network images.} We corrected both counts to the 185 released network maps. They are the atlas maps and the ICA components.
\revchanges{Methods, Datasets, \msl{132}{144}. Results, Contrastive Model: Network Labeling, \msl{785}{797}, and Autoencoder, \msl{876}{880}. Figs.~\msfign{4} and \msfign{5} captions (\mspp{15}{17}). Supplementary Information, NeuroVault preprocessing, \msl{1229}{1246} and \msl{1262}{1264}, as for Reviewer 1's comment \cmt{1.5}.}

\revcomment{2.8}{Fig.~5 caption}\begin{revquote}
2.The caption of Fig. 5 reports both 0.78 and 0.85 for the NeuroVault AUC.
\end{revquote}\response The two values were the two retrieval directions, but the caption did not say so. The corrected caption now names the curves. With each paper's best-matching map, the AUC is 0.85 from text-to-brain and 0.79 from brain-to-text (Fig.~\ref{fig:fig5}b). The brain-to-text value changed from 0.78 because p-maps are inverted and the model was retrained (\cmt{1.3}). Each curve has its own dashed chance curve, from texts shuffled across papers.

\textbf{Lower bound.} The lower bound keeps the map that the model scores lowest for each paper, so that map ranks below a typical candidate even when the texts are shuffled. The lower bound and its chance curve therefore lie below 0.50, and the model beats this chance level too.
\rfig{fig:fig5}{figures/2.8_fig5_perf_quant.png}{\textbf{New \msfig{5}.} Contrastive recall on (a)~PubMed and (b)~NeuroVault, (c)~autoencoder reconstruction, and (d)~network classification. In (b), each paper's target is its best map (upper bound), every map (multi-target), or its worst map (lower bound). Dashed lines show each rule's chance, the same rule after each paper's text is given to another paper, as in Fig.~\ref{fig:bounds}.}
\revchanges{\msfig{5} caption (\msp{17}). Results, Contrastive Model: NeuroVault, \msl{803}{816}, and PubMed, \msl{821}{830}.}

\revcomment{2.9}{SPECTER naming}\begin{revquote}
3.The use of “SPECTER,” “SPECTER2,” and “Specter” is inconsistent. The exact model used in each component should be clarified.
\end{revquote}\response Every text-encoding step uses SPECTER2, the allenai/specter2\_aug2023refresh model with its ad hoc query adapter, and the original SPECTER is never used. We corrected every mention to SPECTER2. The other language models are NeuroQwen3-0.6B, which is Qwen3-0.6B adapted to PubMed titles and abstracts and writes the captions (\cmt{2.11}), and Qwen3.5-27B, which wrote the paper summaries and the synthetic and curated texts (\cmt{2.6}). Reviewer 1 raised the same concern (\cmt{1.6}).
\revchanges{Introduction, \msl{107}{111}. Methods, Text Encoder, \msl{220}{227}, and Text-to-Brain, \msl{364}{373}. Supplementary Information, Text preprocessing, \msl{1131}{1134}. \msfig{S1} caption (\msp{24}), as for Reviewer 1's comment \cmt{1.6}.}

\revcomment{2.10}{LoRA analogy}\begin{revquote}
4.In Section 2.9.2, describing an MLP projection head as similar to a LoRA adapter is technically inaccurate.
\end{revquote}\response We agree and deleted the LoRA analogy. The head is a small network trained on frozen SPECTER2 embeddings, and the revision calls it the text-to-brain head rather than an adapter. Reviewer 1 raised the same concern (\cmt{1.7}).
\revchanges{Methods, Text-to-Brain, \msl{364}{385} (LoRA sentence deleted, projection head named). Supplementary Information, Projection Head, \msl{1328}{1347}, as for Reviewer 1's comment \cmt{1.7}.}

\revcomment{2.11}{Training details}\begin{revquote}
5.The authors should report the model parameters, GPU hardware, training time, random seeds, Qwen fine-tuning strategy, and relevant training hyperparameters.
\end{revquote}\response Table~\ref{tab:training}, the new Supplementary \mstab{S7}, gives the parameters, training data, epochs, optimizer, learning rate, batch size, seed, and training time of every component. All components except the language model train on one RTX 5090 GPU, with fixed seeds and deterministic kernels. The whole pipeline, five-fold cross-validation included, trains from scratch in about 1.8 hours.

\textbf{Qwen fine-tuning strategy.} NeuroQwen3-0.6B is Qwen3-0.6B with every weight fine-tuned, without adapters, on PubMed titles and abstracts for one epoch, which took 14 hours on one H100 GPU (Table~\ref{tab:training}). Its training script removes the validation, test, and NeuroVault evaluation papers from the corpus. Our response to Reviewer 1 reports a memorization check (\cmt{1.2}).
% source: z47 fact sheet (scripts/outputs/runs/final/*.summary.json, scripts/models/*.py defaults); NeuroQwen3-0.6B row: z53 (01_dapt.ipynb, model card)
\rtab{\textbf{Training details.} Rows give the trainable parameters, training data, epochs, optimizer settings, seed, and training time of each component of the final model. Times are for one RTX 5090 GPU, except for NeuroQwen3-0.6B.\label{tab:training}}{\scriptsize\begin{tabular}{L{2.4cm}L{1.75cm}L{2.5cm}L{2.4cm}L{2.3cm}L{0.9cm}L{1.1cm}}\toprule
\textbf{Component} & \textbf{Parameters} & \textbf{Training data (train / validation)} & \textbf{Epochs (selection)} & \textbf{Optimizer, learning rate, batch} & \textbf{Seed} & \textbf{Time} \\\midrule
Autoencoder & 59.9 M & PubMed maps, 23,894 / 2,987 & 101 (lowest validation loss) & AdamW, 5e-5, 256 & 0 & 10 min \\
Contrastive heads, text and image & 0.59 M and 0.30 M & text and map latents, 23,894 / 2,987 & 301 (last) & AdamW, 1e-5, 2,048 & 123 & 6 s \\
Text-to-brain head & 0.59 M & text and map latents, 23,894 / 2,961 & 30 of at most 80 (best validation Dice) & AdamW, 1e-3, 512 & 0 & 22 s \\
Anatomically centered head & 3.8 M & synthetic pairs, 6,309 / 64, and PubMed training pairs & 200 of 2,000 (best PubMed validation recall) & AdamW, 2e-4, 2,048 & 123 & 103 s \\
Stage 1 projector (MLP) & 28.3 M & maps and LLM summaries, 22,547 / 2,837 & 4 of at most 30 (lowest validation loss, early stopping) & AdamW, 3e-4 with cosine decay, 32 & 17 & 17 min \\
Stage 2 projector (MLP) & 28.3 M & concept texts and maps, 972 / 134 & 30 (last) & AdamW, 1e-4 with cosine decay, 18 & 3 & 5 min \\
NeuroQwen3-0.6B & 596 M & PubMed titles and abstracts, 1.22 M / 12,317 (385 M tokens) & 1 & AdamW, 5e-6 with cosine decay, 16 sequences of 1,024 tokens, bf16 & 12345 & 14 h (one H100) \\
\bottomrule\end{tabular}}
\revchanges{Methods, Training, \msl{174}{190} (exclusions, the evaluation papers removed from the language-model corpus, and training time). Methods, Text Decoder, \msl{274}{283} (NeuroQwen3-0.6B training). Supplementary Information, Autoencoder training, \msl{1319}{1322}. New Supplementary \mstab{S7} with the training details per component (\msp{42}).}

\revcomment{2.12}{Random results and failure cases}\begin{revquote}
6.Figs. 2 and 3 mainly show selected qualitative examples. Randomly selected results and representative failure cases should also be provided.
\end{revquote}\response We added random results and failure cases for every evaluation set. In the manuscript, \msfig{3} draws its rows by a fixed rule from the high, medium, and low thirds of the similarity distribution, and its Low rows are failure cases (\cmt{3.9}). Fig.~\ref{fig:failures} shows the six worst brain-to-text cases on the PubMed test set. Random PubMed test papers, the worst text-to-brain cases, the weakest atlas networks, and NeuroVault papers that fail on their best-matching map are in the code repository [author: or name the supplementary figure that holds them]. Failures are common, with 16\% of PubMed test-set captions scoring a negative similarity to the paper's text.
% the other five sets: scripts/figures/review/reviewer_02/2.12_{pubmed_random,pubmed_t2b_worst,networks_bottom,neurovault_b2t_best,neurovault_t2b_best}.png

\textbf{Data and model failures.} Many failures start in the data. For 37\% of the failing maps, the paper's own text and map are among the least similar pairs of the PubMed test set, against 5\% of the successful maps, and Neurosynth, which shares nothing with NeuroVLM, also predicts these maps poorly from their text. Maps with few or scattered foci fail no more often than others. The remaining failures are model misses.
% source: core/a7_failure_cases_percentiles/failure_decomposition.py (outputs/failure_decomposition.md): mismatch = true-pair
% cosine of title + abstract with the paper's map below the test-set 10th percentile; 180 of 482 failures (37.3%) against
% 22 of 482 random non-failures (4.6%); poor input map 28.8% against 28.6%; Neurosynth r below its 10th percentile 31.3% against 5.0%

\textbf{Invented names.} The failure cases showed invented network and region names in the titles. Title decoding is steered toward canonical names. At the title position, tokens that cannot continue a canonical network or region name are masked, and the model's own probabilities choose among the rest.

\textbf{Body text.} The body text and {[}FUNCTION{]} titles are generated without this mask and can contain invented words or stray fragments, such as ``useary'' in Fig.~\ref{fig:failures}. The Discussion states this limitation.

\textbf{Non-Latin text.} A few captions contained non-Latin characters, which come from the language model's multilingual vocabulary. Decoding is restricted to Latin-script tokens.

\revchanges{Figs.~\msfign{2}, \msfign{3}, and \msfign{S4} (\msp{11}, \msp{12}, and \msp{35}), with the Low rows of \msfig{3} as failure cases. Methods, Text Decoder, \msl{298}{305} (title decoding steered toward the canonical names, and Latin-script decoding). Supplementary Information, Brain-to-Text Projector, \msl{1457}{1464}. Discussion, limitations, \msl{949}{957} (failure rate, failures that start in the data, typical failures, and invented words or stray fragments in the body text).}
% the changes line comes before this multi-page figure, so that it is not left alone on a page after the figure pages
\rfigparts{fig:failures}{figures/2.12_pubmed_b2t_worst}{\textbf{Worst brain-to-text cases.} The six PubMed test maps whose caption is least similar to the paper's text, two per page. Each shows the map in sagittal, coronal, and axial views, the paper's title and summary beside the complete generated caption, and the caption's similarity with its percentile.}

\reviewer{3}
\revcomment{3.0}{Summary}\begin{revquote}
Hammonds et al. introduce a meta-analysis framework for aligning text-neuroimage pairs in a shared latent space. The framework can retrieve one modality from the other by comparing and ranking the latent vector embeddings. The framework can also generate one modality from the other by passing the latent vector embeddings through generative text and brain decoders. The authors tested the framework on multiple neuroimaging datasets. They found statistically significant retrieval and generation performance with sensible qualitative examples.
\end{revquote}\response This summary is accurate. Our responses to the comments follow (\cmt{3.1} to \cmt{3.10}).

\revcomment{3.1}{Novelty and what the model learns}\begin{revquote}
Augmenting existing meta-analysis tools with the flexibility provided by LLMs is an important and promising research direction. The manuscript appears technically sound and the NeuroVLM architecture seems appropriate for integrating the contrastive and generative models. However it is less clear how NeuroVLM compares to previous meta-analysis methods. The manuscript would benefit from a) clarifying the novelty of the framework by demonstrating the importance of generation and context and b) interpreting what the model learns.
\end{revquote}\response Generation and context both matter, and the model learns interpretable associations between terms and brain maps.

\textbf{1. Generation and context.} Generation matters. Generated maps match each test paper's map better than the nearest training paper's map does (Reviewer 1, \cmt{1.1}). Generated captions are as close to the paper's text as retrieved captions and rarely report another study's details, which retrieved captions often do (\cmt{3.3}). Context matters too. SPECTER2 beats a context-free word2vec encoder on every measure, by a modest margin (\cmt{3.4}).

\textbf{2. What the model learns.} The shared latent dimensions are interpretable, since each clear dimension decodes to a map that matches the Neurosynth map of its top term (\cmt{3.5}). Swapping one word in a text moves the generated map toward the matching Neurosynth contrast, but negation has no effect, so the model composes terms additively. Against previous meta-analysis methods, NeuroVLM gives closer text-to-brain maps than Neurosynth and NeuroQuery (\cmt{3.2}).
\revchanges{Introduction, \msl{99}{105}, Results, Generative Models: PubMed, \msl{606}{613}, and Discussion, \msl{913}{918}, which state the contribution of generation and context.}

\revcomment{3.2}{Comparison with existing methods}\begin{revquote}
My main concern is about the scope of the manuscript. NeuroVLM is clearly a useful tool but it is difficult to evaluate how much it improves on previous work. If the main contribution is higher performance then it would be helpful to include more direct comparisons. The only comparison between NeuroVLM and existing methods is in Figure S3 where it had comparable performance to NiCLIP and NeuroConText (albeit with fewer parameters). If the main contribution is instead more capabilities then it would be helpful to emphasize what NeuroVLM can do that previous methods cannot.
\end{revquote}\response NeuroVLM improves on existing methods in both respects, with higher performance and more capabilities.

\textbf{1. Direct comparisons.} Each method turns a test paper's title and abstract into a map, which is scored against that paper's map (Fig.~\ref{fig:baselines}). NeuroVLM gives the closest maps, ahead of Text2Brain, NeuroQuery, and Neurosynth (Dice 0.53 against at most 0.49). In retrieval, NeuroVLM is slightly ahead of NiCLIP and NeuroConText on the test papers that neither was trained on (Fig.~\ref{fig:figS3}).

\textbf{2. New capabilities.} We know of no other model that writes text from group-level brain maps. NeuroVLM also retrieves and generates in both directions through one shared space, and its generated maps and captions match or beat retrieved ones (\cmt{3.3}). The Introduction states these capabilities.
\rfig{fig:figS3}{figures/3.2_figS3_recall_comparison.png}{\textbf{New \msfig{S3}.} Recall@k for NeuroVLM (blue), NiCLIP (orange), and NeuroConText (green) on the 366 PubMed test papers that none of the three models was trained on. Solid lines are text to image and dashed lines image to text.}
\revchanges{Introduction, \msl{99}{105}, on the new capabilities. Results, Generative Models: PubMed, \msl{606}{613}, and Contrastive Model: PubMed, \msl{821}{830}. Supplementary Information, Model Comparison, with \msfig{S3} redrawn from the final model (\msp{34}). New Supplementary Section Text-to-Brain Baselines and Retrieval, \msl{1935}{1953}, with \msfig{S6} (\msp{44}).}

\revcomment{3.3}{Generation versus retrieval and granularity}\begin{revquote}
One capability that sets NeuroVLM apart from previous methods is that it can generate original text and neuroimages rather than only retrieving from a fixed set. However it would be helpful to show that the generated outputs are in fact more accurate and useful than retrieved outputs. For example if NeuroVLM uses an overfit decoder that collapses into a fixed set of outputs then it might be no better than existing retrieval methods. In fact the canonical projection matrix might explicitly drive the framework towards more retrieval-like outputs. Is there a way to quantify the granularity of NeuroVLM based on its output space?
\end{revquote}\response Generated outputs are at least as accurate as retrieved ones. The decoder does not collapse, although the canonical projection pulls about one caption in ten toward a prototype, and we quantify the granularity of the outputs below. Reviewer 1 raised the same comparison (\cmt{1.1}).

\textbf{1. Generated versus retrieved outputs.} Generated maps are closer to each test paper's map than the retrieved map, the nearest training paper's map (Table~\ref{tab:t2b}). By MiniLM cosine, generated captions are at least as close to the paper's text as the nearest training caption (Table~\ref{tab:genret} and Fig.~\ref{fig:granularity}b). BERTScore favors training captions, even random ones, for their summary style. Retrieved captions often carry another study's details, and generated captions rarely do (57\% against 2\%).

\textbf{2. Collapse.} The decoder does not collapse to a fixed set of outputs. Over 99\% of stage~2 test captions are unique, against 73\% of retrieved captions (Table~\ref{tab:genret}). The language model generates every caption, and only the network and region titles are steered toward canonical names, by masking tokens that cannot continue one (Reviewer 2, \cmt{2.12}). The body text is not masked.

\textbf{3. Canonical projection.} The canonical projection pulls captions toward its prototypes. With it, about one caption in ten nearly copies a curated training text, against almost none without it (Fig.~\ref{fig:granularity}c). This is the price of the drop in study details under point 4, and the Methods state the trade-off.

\textbf{4. Granularity.} Yes. We quantify granularity by the spread of titles and by the study detail in the body text. Test captions use nearly every canonical network and region name, and no single title is given to more than one test map in ten. The share of captions that name a study detail, such as a sample size, patient group, or task, falls from 36\% at stage~1 to 2\% at stage~2 (Fig.~\ref{fig:granularity}a). Stage~2 captions describe the map, not a single study.
\rfig{fig:granularity}{figures/3.3_granularity_tradeoff.png}{\textbf{Granularity of generated captions.} Captions of PubMed test maps and NeuroVault papers, retrieved or generated at each stage. Panels show (a)~the share naming a study detail, and the MiniLM cosine (b)~with the paper's own text and (c)~with the nearest training text.}
% source: core/a2_granularity_tradeoff/outputs/metrics.json and table.csv (final model; PubMed test set, n = 2,971;
% Unique = unique_frac_text: stage 2 0.994, stage 1 0.969, retrieved 0.733);
% random training text: core/a1_brain_to_text_baselines/outputs/brain_to_text_baselines_table.md (model-independent)
\rtab{\textbf{Generated versus retrieved captions.} Texts returned for PubMed test maps, scored against each paper's title and LLM summary by MiniLM cosine and BERTScore F1. Brackets give the MiniLM cosine against another paper's text. Unique is the share of distinct texts.\label{tab:genret}}{\small\begin{tabular}{lccc}\toprule
\textbf{Text returned for a test map} & \textbf{MiniLM cosine (null)} & \textbf{BERTScore F1} & \textbf{Unique} \\\midrule
NeuroVLM, stage 2 caption (paper) & 0.463 (0.397) & 0.553 & 99\% \\
NeuroVLM, stage 1 caption & 0.431 (0.358) & 0.577 & 97\% \\
Retrieval, training caption nearest to the map & 0.428 (0.362) & 0.588 & 73\% \\
Random training text & 0.363 & 0.580 & -- \\
\bottomrule\end{tabular}}
\revchanges{Methods, Text Decoder, \msl{315}{319}, on the canonical projection, which removes study details but adds near-copies. Discussion, \msl{913}{918}. New Supplementary Section Caption Granularity, \msl{2165}{2174}, with \msfig{S8} (\msp{48}).}

\revcomment{3.4}{Context-free baseline}\begin{revquote}
Another capability that sets NeuroVLM apart is the context and flexibility provided by the LLM encoder. However it would be helpful to directly quantify how the LLM encoder performs relative to context-free baselines (such as mean Word2Vec embeddings) on the retrieval and text-to-neuroimage tasks.
\end{revquote}\response SPECTER2 outperforms a context-free baseline on both tasks, by a modest margin (Fig.~\ref{fig:encoder}). We replaced SPECTER2 with the mean word2vec embedding of each title and abstract and retrained only the projection heads. Mean retrieval AUC, the area under the recall curve, is 0.83 against 0.81. Text-to-brain Dice and voxel correlation are also higher with SPECTER2. We added this comparison to the Discussion and the Supplementary Information.
\rfig{fig:encoder}{figures/3.4_encoder_swap.png}{\textbf{Context-free encoder.} PubMed test-set performance with SPECTER2 or with the mean word2vec embedding of each title and abstract, after retraining only the projection heads. Panels show (a)~retrieval AUC in both directions and (b)~text-to-brain Dice and voxel correlation.}
\revchanges{Introduction, \msl{104}{105}, and Discussion, \msl{900}{906}. New Supplementary Section Context-Free Encoder, \msl{2198}{2204}, with \msfig{S9} (\msp{49}).}

\revcomment{3.5}{Interpretation, perturbations, and ablations}\begin{revquote}
It would be helpful to include more analyses like Figure 6 that interpret different components of the model. Interpreting the encoder weights could reveal the specific associations between text and neuroimage elements that NeuroVLM learned. Perturbing the inputs (such as adding negations) could show how NeuroVLM goes beyond linear additivity. Providing ablation results in the supplementary materials could compare the importance of different modeling decisions (such as the cross-entropy loss and canonical projection matrix).
\end{revquote}\response The shared latent dimensions are interpretable, word swaps move the generated map but negation does not, and the cross-entropy loss matters most for retrieval.

\textbf{1. Encoder weights.} The encoder weights act on SPECTER2 features that have no fixed meaning, so we interpret the latent dimensions that texts and maps share instead. For the clearest dimensions, the decoded map matches the Neurosynth map of the dimension's top term, such as listening in superior temporal cortex (Fig.~\ref{fig:latent}, $r = 0.65$). The text and brain sides agree on what each dimension means. This figure is new in the Supplementary Information.

\textbf{2. Perturbations.} Word swaps and joined concepts move the map as expected, but negation does not (Fig.~\ref{fig:probes}). Swapping one word moves the map toward the matching Neurosynth contrast, and joining two concepts covers both. A sentence and its negation give nearly the same map ($r = 0.84$ to 1.00). NeuroVLM therefore does not go beyond linear additivity, and the Discussion states this limitation.

\textbf{3. Ablations.} Table~\ref{tab:ablation}, new in the Supplementary Information, ablates the losses. The cross-entropy loss matters most for retrieval, since retrieval AUC drops from 0.83 to 0.62 without it, while the temperature and a one-directional loss barely matter. The final text-to-brain head beats heads trained with the other losses. Our response to \cmt{3.3} gives the effect of the canonical projection on captions.
\rfig{fig:latent}{figures/3.5_latent_interpretation.png}{\textbf{Latent dimensions.} (a)~Decoded maps of six latent dimensions, each labeled with its three most associated Cognitive Atlas terms and its correlation $r$ with the top term's Neurosynth map. (b)~Agreement of each dimension's term associations between the text and brain sides, against shuffled dimensions.}
\rfig{fig:probes}{figures/3.5_perturbation_probes.png}{\textbf{Perturbation probes.} (a)~NeuroVLM maps for ``auditory stimulation'' and ``visual stimulation'', each averaged over eight sentences that differ only in this phrase, their difference, and the difference of the Neurosynth ``auditory'' and ``visual'' term maps. One axial slice shows each map's top 5\% of voxels, and $r$ is the correlation of the two differences, averaged over the sentences. (b)~This correlation for five word pairs, against the other pairs' Neurosynth differences. (c)~Dice overlap of each ``A and B'' map with the union of the A and B maps, against the better of A or B alone.}
% source: core/a5_loss_ablation/outputs/a5_loss_ablation_table.md and a5_loss_ablation_text_to_brain.csv (3 seeds, test set, heads retrained on the final latents).
% The final text-to-brain head also differs from the other three in optimizer settings, epochs and model selection (hence the row label).
\rtab{\textbf{Loss ablations.} Heads retrained on the final model's latents with the loss shown and scored on the PubMed test set. Values are means of three seeds.\label{tab:ablation}}{\small\begin{tabular}{lc}\toprule
\textbf{Loss} & \textbf{PubMed test set} \\\midrule
\multicolumn{2}{l}{\textbf{Contrastive heads, retrieval AUC}} \\
\quad Symmetric InfoNCE cross-entropy, $\tau = 0.07$ (final) & 0.833 \\
\quad One direction only, text $\rightarrow$ brain / brain $\rightarrow$ text & 0.834 / 0.827 \\
\quad $\tau = 0.03$ / $\tau = 0.20$ / learnable $\tau$ & 0.829 / 0.827 / 0.833 \\
\quad No cross-entropy, cosine of matched pairs only & 0.618 \\\midrule
\multicolumn{2}{l}{\textbf{Text-to-brain head, Dice}} \\
\quad Latent MSE and map correlation, final head with its own training settings & 0.530 \\
\quad Latent MSE only & 0.512 \\
\quad Voxel-wise binary cross-entropy through the decoder & 0.491 \\
\quad Latent cosine only & 0.385 \\
\bottomrule\end{tabular}}
\revchanges{Introduction, \msl{76}{79}, where the contrast with additive term models no longer implies non-linear composition. Discussion, limitations, \msl{956}{961}, with a new sentence that NeuroVLM ignores negation. New Supplementary Section Latent Dimensions and Loss Ablations, \msl{2231}{2240}, with \msfig{S10} (\msp{50}) and \mstab{S10} (\msp{51}).}

\revcomment{3.6}{Network, anatomy, and cognition outputs}\begin{revquote}
It is interesting that the model can generate separate brain-to-text outputs corresponding to the predicted network, anatomy, and cognition. How does quantitative performance compare between these outputs? What if the text-to-brain outputs were generated using only text about the network, anatomy, or cognition? These analyses would also shed light on the associations learned by NeuroVLM.
\end{revquote}\response {[}REGION{]} and {[}FUNCTION{]} titles are correct more often than {[}NETWORK{]} titles, and maps generated from one kind of text alone beat a null baseline. Every title is generated by the language model, with {[}NETWORK{]} and {[}REGION{]} titles steered toward canonical names by masking tokens that cannot continue one (Reviewer 2, \cmt{2.12}).

\textbf{1. Comparison of the three outputs.} A {[}NETWORK{]} title is correct if it names the map's network, a {[}REGION{]} title if it names a region inside the map, and a {[}FUNCTION{]} title if it names a function that [author: how a function is decoded from the map, e.g. "Neurosynth decoding of the map ranks in its top terms"] (Fig.~\ref{fig:steering}). On PubMed test maps, {[}REGION{]} titles are correct most often (63\% against 24\% by chance) and {[}NETWORK{]} titles least. Fig.~\ref{fig:steering} also scores atlas networks and ICA components.

\textbf{2. Network, anatomy, or cognition text only.} Maps generated from a network name, a region name, or a cognitive term alone all match their references better than a null baseline (Table~\ref{tab:onekind}). Region maps overlap their references least, because most regions are smaller than the top 10\% of voxels that Dice compares.
\rfig{fig:steering}{figures/3.6_steering_correctness.png}{\textbf{Correctness of generated titles.} Share of {[}NETWORK{]}, {[}REGION{]}, and {[}FUNCTION{]} titles judged correct for atlas networks, ICA components, and PubMed test maps. Dashed lines show chance, from titles shuffled across maps.}
% source: core/a6_modality_restricted_t2b/outputs/modality_table.md (final text-to-brain head, calibrated maps)
\rtab{\textbf{Text-to-brain from one kind of text.} Maps generated from bare network names, region names, or Cognitive Atlas terms, scored against the matching atlas map, region, or Neurosynth term map. Brackets give a null baseline, the same maps scored against the other items' references.\label{tab:onekind}}{\small\begin{tabular}{llccc}\toprule
\textbf{Input text} & \textbf{Reference} & \textbf{$n$} & \textbf{Dice (null)} & \textbf{Pearson $r$ (null)} \\\midrule
Network name & Yeo-7 and network atlas maps & 126 & 0.29 (0.05) & 0.29 ($-$0.01) \\
Region name & Harvard-Oxford and AAL regions & 178 & 0.07 (0.01) & 0.15 (0.00) \\
Cognitive Atlas term & Neurosynth map of the term & 130 & 0.28 (0.13) & 0.25 (0.02) \\
\bottomrule\end{tabular}}
\revchanges{Methods, Text Decoder, \msl{298}{305}, on title decoding steered toward the canonical names. Discussion, limitations, \msl{949}{957}, where {[}NETWORK{]} titles are the least reliable, especially for single-study maps.}

\revcomment{3.7}{Unsupported claims}\begin{revquote}
Some claims are not justified in the context of the study. For example “each component of NeuroVLM is modular and interchangeable with larger or smaller models depending on requirements” (line 101) and “our results underscore that the limiting factor for biological fidelity is often the structure of what the field reports and shares” (line 672). While intuitively true these statements would be better supported by including scaling and ablation analyses.
\end{revquote}\response We narrowed both claims to what the study supports.

\textbf{1. Modular and interchangeable components.} The sentence claims only that the text encoder can be swapped by retraining the small projection heads. The encoder swap of \cmt{3.4} shows this, where replacing SPECTER2 with word2vec cost a modest drop in performance (Fig.~\ref{fig:encoder}). The reference to larger or smaller models is removed, because we tested no scaling.

\textbf{2. Limiting factor.} The sentence was meant to contrast coordinate-based supervision with full statistical maps. It is rewritten to say exactly that, with citations on what coordinate-based meta-analysis loses, and no longer claims a general limiting factor.
\revchanges{Introduction, \msl{109}{111}, modularity sentence. Discussion, \msl{990}{998}, limiting-factor sentence, and \msl{965}{967}, where the sentence on our strongest limitations is deleted.}

\revcomment{3.8}{Titles and abstracts only}\begin{revquote}
Could the authors justify why NeuroVLM was only trained on titles and abstracts? I imagine that other sources of information (such as table captions) might be useful for checking whether the activation coordinates are actually related to the claims in the titles and abstracts.
\end{revquote}\response Titles and abstracts exist for every paper, but table captions exist only for the papers with open-access full text, fewer than half. Captions also say little about where the coordinates lie, so they cannot check the abstract's claims (Table~\ref{tab:captions}). A classifier that predicts a table's main region from its caption performs at chance (macro-F1 0.06 against 0.05), while the same classifier does better from the abstract. The Discussion gives these reasons.
\rtab{\textbf{Uninformative table captions.} Captions of coordinate tables from open-access papers, picked by hand from 60 random examples. The right column lists the Harvard-Oxford regions that hold most of each table's peaks.\label{tab:captions}}{\small\begin{tabular}{lp{6.6cm}p{5.2cm}}\toprule
\textbf{Kind} & \textbf{Caption} & \textbf{Where its coordinates lie} \\\midrule
boilerplate & Brain regions revealed by contrasts of interest. & anterior cingulate, frontal pole, insula \\
boilerplate & The main results of group analysis of block data. & posterior cingulate, amygdala, thalamus \\
boilerplate & Regional imaging relationships. & lateral occipital cortex, precuneus \\
condition code & Simple Effects contrast JUDGE –BODY. & anterior cingulate, lateral occipital, postcentral \\
condition code & Neural activations for the contrast AU $>$ OC. & lateral occipital, paracingulate, posterior cingulate \\
statistics only & Functional results (p = .05 family-wise error (FWE) cluster corrected). & lateral occipital, pallidum \\
statistics only & Significant clusters exhibiting interactions of interest at p $<$ 0.001 (uncorrected), n = 15. & insula, paracingulate, orbitofrontal \\
region list only & Significant clusters from amygdala connectivity analyses. & medial frontal cortex, insula \\
\bottomrule\end{tabular}}
\revchanges{Discussion, \msl{927}{934}, on table captions and on the pooling of a paper's tables as a limitation.}

\revcomment{3.9}{Percentiles of the qualitative examples}\begin{revquote}
It would be helpful to know where the qualitative examples in Figures 2 and 3 fall within the distribution of cosine similarity quality scores. Alternatively the authors could provide a percentile for each example.
\end{revquote}\response Table~\ref{tab:percentiles} gives each example's percentile within its set, from the 16th to the 99th, and \msfig{3} plots both similarity distributions with each example marked. The rows of \msfig{3} follow a fixed rule. Each third of the distribution contributes the papers nearest its center that have a clean generated text, and the High and Medium rows also require a title that names one domain and a true map with distinct clusters.
% source: core/a7_failure_cases_percentiles/outputs/a7_example_percentiles.md (final model)
\rtab{\textbf{Percentiles of the examples.} Percentile of each example of \msfig{3} and \msfig{2} among all items of its set, by text-to-brain and brain-to-text cosine similarity. \msfig{3} ranks 305 NeuroVault papers and \msfig{2} ranks atlas networks.\label{tab:percentiles}}{\small\begin{tabular}{lcc}\toprule
\textbf{Example} & \textbf{Text $\rightarrow$ brain percentile} & \textbf{Brain $\rightarrow$ text percentile} \\\midrule
\msfig{3}, High, working memory & 86 & 79 \\
\msfig{3}, High, language & 77 & 87 \\
\msfig{3}, Medium, memory & 52 & 62 \\
\msfig{3}, Medium, reward & 36 & 43 \\
\msfig{3}, Low, poor input map & 16 & 18 \\
\msfig{3}, Low, model miss & 22 & 19 \\
\msfig{2}, Visual & 76 & 99 \\
\msfig{2}, Language & 75 & 37 \\
\msfig{2}, Auditory & 31 & 19 \\
\msfig{2}, Somatomotor & 39 & 72 \\
\msfig{2}, Interoception & 44 & 55 \\
\msfig{2}, Default mode & 69 & 94 \\
\msfig{2}, Frontoparietal control & 48 & 53 \\
\msfig{2}, Salience & 33 & 46 \\
\bottomrule\end{tabular}}
\revchanges{\msfig{3}, on one page, with rows chosen by a fixed rule and distribution strips that mark each example (\msp{12}). \msfig{S4} caption, which gives the percentiles of its components (\msp{35}). Results, Generative Models: Networks, \msl{446}{448}, with the percentiles of the \msfig{2} examples, and NeuroVault, \msl{458}{460}, on the complete captions and the choice of rows.}

\revcomment{3.10}{Code availability}\begin{revquote}
I was able to install and run the code. The project is very well-documented and user-friendly. I believe that it provides a very usable resource for the community.
\end{revquote}\response We also made the reproducibility changes described in our response to Reviewer 1 (\cmt{1.11}).

\FloatBarrier
\end{document}
